% Generated draft: requires mathematical and source review.
\documentclass[11pt]{article}
\usepackage[margin=25mm]{geometry}
\usepackage{fontspec}
\setmainfont{Latin Modern Roman}
\usepackage{amsmath,amssymb,mathtools,unicode-math}
\setmathfont{Latin Modern Math}
\usepackage{enumitem,needspace,xurl,hyperref}
\setlength{\parindent}{0pt}
\setlength{\parskip}{5pt}
\setlength{\emergencystretch}{3em}
\newcommand{\N}{\mathbb{N}}
\newcommand{\Z}{\mathbb{Z}}
\newcommand{\Q}{\mathbb{Q}}
\newcommand{\R}{\mathbb{R}}
\newcommand{\C}{\mathbb{C}}
\newcommand{\F}{\mathbb{F}}
\newcommand{\A}{\mathbb{A}}
\newcommand{\PP}{\mathbb{P}}
\newcommand{\E}{\mathbb{E}}
\newcommand{\eps}{\varepsilon}
\newcommand{\dd}{\,\mathrm{d}}
\begin{document}
\section*{UnsolvedMath ID: 30004036}
\textbf{Source problem code:} catalogue entry
\subsection{Definitions and mathematical statement}
The supplied record does not explicitly define its graph conventions, \(q\)-equivalence, the operation \(+\), or unqualified Ramsey equivalence. The following is a conditional formalization under the usual edge-Ramsey interpretation; these conventions require confirmation against the cited sources.

Take graphs to be finite, simple, and undirected. For graphs \(G,F\) and an integer \(q\geq2\), write
\[
G\longrightarrow(F)_q
\]
if every map \(c:E(G)\to\{1,\ldots,q\}\) admits a subgraph of \(G\) isomorphic to \(F\) whose edges all receive the same color. Copies are not required to be induced, and the coloring need not use every color. Define
\[
F\sim_q F'
\quad\Longleftrightarrow\quad
\text{for every finite simple undirected graph }G,
\quad
G\longrightarrow(F)_q\ \Longleftrightarrow\ G\longrightarrow(F')_q.
\]
Under this interpretation, \(q\)-equivalent means related by \(\sim_q\). A graph is connected if it has at least one vertex and every two vertices are joined by a path. Write \(F\not\cong F'\) when the graphs are not isomorphic. Let \(K_k\) be the complete graph on \(k\) vertices. Interpret \(K_k+H\) as the disjoint union of \(K_k\) and \(H\), with no edges between the two components, and interpret unqualified Ramsey equivalence as \(\sim_2\).

With these conventions, the supplied questions are:
\begin{enumerate}
\item Is it true that, for every pair of connected graphs \(H,H'\),
\[
H\sim_3 H'\quad\Longrightarrow\quad H\sim_2 H'?
\]
\item Is it true that
\[
\forall q\in\Z,\ q\geq2,\quad
\exists\text{ connected graphs }H,H'
\quad\text{such that}\quad
H\not\cong H'\ \text{ and }\ H\sim_q H'?
\]
The graphs may depend on \(q\); the wording does not require one pair to work for every \(q\).
\item For a fixed parameter \(k\), characterize the graphs \(H\) for which
\[
K_k\sim_2(K_k+H),
\]
or, equivalently, for which every finite simple undirected graph \(G\) satisfies
\[
G\longrightarrow(K_k)_2
\quad\Longleftrightarrow\quad
G\longrightarrow(K_k+H)_2.
\]
\end{enumerate}
The last question does not explicitly restrict \(H\) to connected graphs. The source supplies neither the intended range of \(k\) nor a restriction excluding the empty graph as \(H\); these boundary conditions remain unresolved. The notation \(K_k\) ordinarily requires a nonnegative integer vertex count, but no intended lower bound is inferred here. If \(+\) denotes a join, or if unqualified Ramsey equivalence refers to a different color quantifier, the third question must be reformulated accordingly.
\subsection{Short English statement}
Under the explicitly provisional edge-Ramsey conventions, ask whether connected graphs equivalent for three colors must be equivalent for two, whether each integer \(q\geq2\) admits nonisomorphic connected \(q\)-equivalent graphs, and which graphs can be added disjointly to \(K_k\) without changing its two-color Ramsey equivalence class. The intended graph conventions, meaning of \(+\), color convention in the final question, and range of \(k\) require source confirmation. The supplied open status is dated curation, not an independently verified conclusion.
\subsection{Sources}
\begin{itemize}
\item Local registry: \texttt{data/unsolvedmath/problems.json}, supplied record with permanent UnsolvedMath ID \(30004036\), title ``Equivalence Hierarchies for Connected Graphs,'' exact source problem code \texttt{catalogue entry}, in the Oberwolfach Reports Open Problems collection. The record lists proposed year \(2019\) and supplies no individual proposer.
\item Supplied report citation: \emph{Graph Theory (2020)}, Oberwolfach Reports, DOI \texttt{10.4171/owr/2019/1}, pp.\ 5--63. \url{https://doi.org/10.4171/owr/2019/1}. The literature-source entry labels this ``Oberwolfach Report 1/2019.'' The report title's year and the report identifier's year are retained as supplied.
\item Supplied literature reference: Clemens, Liebenau, and Reding, \emph{Ramsey equivalence}. \url{https://arxiv.org/abs/1809.09232}.
\item Literature assessment dated exactly \texttt{2026-08-22}: the supplied record labels the problem open and states that the connected-graph implication and classification of further \(q\)-equivalent graph pairs remain open in the cited literature. This is supplied dated curation, described in the background as literature triage derived from the supplied OWR curation; it has not been independently verified. The record expressly calls for independent verification of time-sensitive status and cited claims.
\item The supplied registry labels the statement \texttt{corrected\_verified} and says its mathematical content was preserved while notation, typography, duplicate clauses, or omitted definitions were repaired for self-contained publication. This is a supplied verification claim, not mathematical verification of this draft or confirmation of the conventions used above.
\end{itemize}
\end{document}
